\section{La superficie de parámetros}
\label{sec:parametros}

Esta sección enumera \emph{todos} los parámetros que existen, con el nodo al que
pertenecen, dónde viven, su tipo y los niveles que se han instanciado de hecho. La
enumeración se obtuvo de la base de datos y no del código, porque lo que decide
qué se ha explorado son las filas y no las firmas de las funciones.

Un parámetro que aparece con un solo nivel instanciado no es un parámetro
explorado: es una constante que nadie eligió, y esa es exactamente la lectura que
la palabra \cod{inherited} le da.

\subsection{Predicción: quince parámetros}

\begin{center}
\small
\begin{longtable}{@{}L{4.9cm} L{1.6cm} L{1.6cm} L{5.6cm}@{}}
\toprule
parámetro & nodo & tipo & niveles instanciados\\
\midrule
\endfirsthead
\toprule
parámetro & nodo & tipo & niveles instanciados\\
\midrule
\endhead
\cod{embedding\_config\_id}      & substrate & uuid    & \textbf{25}. La rejilla completa del sustrato, desglosada en la tabla~\ref{tab:sustrato}\\
\cod{metric}                     & retriever & enum    & 1: \cod{cosine}\\
\cod{limit\_per\_entry}          & retriever & entero  & 1: 200, en los 86 prediction sets. El valor 30 aparece en un único trabajo que falló y no produjo ninguno. Es la profundidad de recuperación, no el corte de evaluación\\
\cod{search\_backend}            & retriever & enum    & 1: \cod{numpy}. El KNN corre siempre en CPU; el valor por omisión \cod{auto} elegía GPU en silencio, que es el único camino que falla\\
\cod{annotation\_set\_id}        & bank      & uuid    & 1: GOA 220. Lo fija el frame\\
\cod{donor\_policy}              & bank      & objeto  & 2: la restrictiva y la vacía. La restrictiva admite once códigos experimentales y de alto rendimiento más \cod{IC} y \cod{TAS}, trece en total; la vacía no filtra por código de evidencia y es por tanto \textbf{la más permisiva de las dos}. El filtro solo se aplica \cod{if codes}, de modo que la lista restringe\\
\cod{exclude\_self\_neighbour}   & retriever & booleano& 2. Retira el auto-acierto a distancia cero\\
\cod{expand\_votes\_to\_ancestors}& retriever& booleano& 2. Propaga el voto de un donante a los ancestros del término\\
\cod{aspect\_separated\_knn}     & retriever & booleano& 2. Recupera vecinos por aspecto en lugar de una sola vez\\
\cod{ontology\_snapshot\_id}     & frame     & uuid    & 1. El grafo bajo el que se resuelven los términos del banco\\
\cod{query\_set\_id}             & frame     & uuid    & 1. El conjunto de queries de la campaña\\
\cod{compute\_alignments}        & features  & booleano& 1: \cod{false} en los 86 prediction sets. El nivel \cod{true} solo aparece en trabajos que no produjeron ninguna fila; la identidad disponible viene del backfill de \cod{sequence\_alignment}\\
\cod{compute\_taxonomy}          & features  & booleano& 1: \cod{false} en los 86 prediction sets, igual que la anterior\\
\cod{compute\_reranker\_features}& features  & booleano& 1: \cod{false} en los 86 prediction sets, igual que las dos anteriores\\
\cod{batch\_size}                & ninguno   & entero  & 1: 1024. Parámetro de ejecución, no de método. No debe entrar en la identidad de un nivel\\
\bottomrule
\caption{Los quince parámetros de \cod{predict\_go\_terms}, con sus niveles
instanciados sobre 69 trabajos.}
\end{longtable}
\end{center}

\begin{aviso}{Las banderas \cod{compute\_} son precondiciones, no ablaciones}
Las tres banderas \cod{compute\_} no cambian la predicción: cambian qué se
\emph{registra} sobre ella. Apagarlas no es una ablación del método, es renunciar
a poder estratificar después. \cod{compute\_alignments} apagado hace ilegibles
cuatro de los siete ejes de estrato, porque los ejes de vecindario necesitan un
donante alineado. La ablación de rasgos, que sí es una decisión de método,
pertenece al nodo \cod{features} y se hace retirando columnas del esquema de
rasgos, nunca apagando el cálculo.
\end{aviso}

\subsection{Evaluación: siete parámetros}

\begin{center}
\small
\begin{tabularx}{\linewidth}{@{}L{5.2cm} L{1.6cm} X@{}}
\toprule
parámetro & nodo & niveles instanciados\\
\midrule
\cod{prediction\_set\_id}      & --       & 46. El arm que se evalúa\\
\cod{evaluation\_set\_id}      & frame    & 2. Las dos variantes de la ventana $220 \rightarrow 227$\\
\cod{information\_accretion\_set\_id} & frame & 1\\
\cod{max\_k\_position}         & retriever & \textbf{7}: 1, 2, 3, 5, 10, 20, 30. Es el corte de evaluación, que trunca una lista ya recuperada\\
\cod{th\_step}                 & frame    & 2: 0.0005 y 0.002, más la ausencia del campo, que toma el valor por omisión 0.01\\
\cod{temporal\_window}         & frame    & 1: la cadena \cod{SELECT\_220\_227}. Etiqueta del arnés\\
\cod{frame}                    & --       & 1: la cadena \cod{internal}. Etiqueta del arnés\\
\bottomrule
\end{tabularx}
\end{center}

\subsection{El sustrato, desglosado}

De los once campos que determinan un substrate, solo siete varían de hecho en el
registro, y de esos siete cuatro solo distinguen la familia de modelos
preentrenados de las tres codificaciones aprendidas. El eje real lo llevan
\cod{model\_name} (12 valores), \cod{model\_backend} (7) y \cod{layer\_indices} (12). Es decir: \textbf{el eje del sustrato, tal
como está instanciado, es modelo por capa}.

\begin{center}
\small
\begin{tabularx}{\linewidth}{@{}L{5.3cm} L{2.8cm} L{2.4cm} X@{}}
\toprule
\cod{model\_name} & \cod{backend} & capas & notas\\
\midrule
\cod{facebook/esm2\_t6\_8M\_UR50D}   & \cod{esm}    & 0 & el más pequeño, usado como suelo\\
\cod{facebook/esm2\_t33\_650M\_UR50D}& \cod{esm}    & 0, 11, 22, 33 & rejilla de capas\\
\cod{facebook/esm2\_t36\_3B\_UR50D}  & \cod{esm}    & 0 & \\
\cod{esmc\_600m}                     & \cod{esm3c}  & 0, 12, 24, 35 & rejilla de capas\\
\cod{ElnaggarLab/ankh-base}          & \cod{ankh}   & 0, 10, 16, 32, 48 & las capas 10, 16 y 32 \textbf{no son evaluables}: puntuación constante\\
\cod{ElnaggarLab/ankh-large}         & \cod{ankh}   & 0 & \\
\cod{Rostlab/prot\_t5\_xl\_half}     & \cod{t5}     & 0, 8, 16, 24 & rejilla de capas\\
\cod{Rostlab/ProstT5}                & \cod{t5}     & 0 & \\
\cod{mila-intel/ProtST-esm1b}        & \cod{protst} & 0 & \\
\cod{rung2-dense}, \cod{rung2-pooled}& \cod{learned-code} & --- & codificaciones aprendidas. \cod{max\_length} 1022, sin normalizar\\
\cod{rung2-residue}                  & \cod{residue-sparse} & --- & idem\\
\bottomrule
\end{tabularx}
\captionof{table}{Las 25 configuraciones de sustrato registradas, en siete
familias de modelo y doce nombres de modelo distintos, nueve preentrenados y tres
aprendidos. \cod{pooling} es \cod{mean} en las veintidós preentrenadas y
\cod{learned} o \cod{residue-sparse-mean} en las tres aprendidas;
\cod{max\_length} es 2048 en las preentrenadas y 1022 en las aprendidas.}
\label{tab:sustrato}
\end{center}

\begin{aviso}{Un confundido que invalida el eje del sustrato tal como está hoy}
Doce de las configuraciones tienen 528{,}294 embeddings almacenados y once tienen
60{,}797, un factor de \textbf{8.69}. Una comparación entre un arm con el banco
grande y otro con el banco pequeño mide dos cosas a la vez, la representación y el
tamaño del banco de donantes, y ninguna cantidad posterior las separa. El eje del
sustrato no puede producir un veredicto hasta que los bancos se igualen. El coste
de igualarlos está en la sección~\ref{sec:presupuesto} y es asumible.
\end{aviso}
